\begin{vidu}
	Cung cấp nhiệt lượng 1,5 J cho một khối khí trong một xilanh đặt nằm ngang. Chất khí nở ra, đẩy pit-tông đi một đoạn 5 cm. Biết lực ma sát giữa pit-tông và xilanh có độ lớn là 20 N, coi pit-tông chuyển động thẳng đều. Tính: 
	\begin{itemize}
		\item[a)] Độ lớn công của khối khí thực hiện.
		\item[b)] Độ biến thiên nội năng của khối khí
	\end{itemize}
	\loigiai{
		\begin{itemize}
			\item[a)] Do pit-tông chuyển động thẳng đều nên lực đẩy $\vec F$ của khối khí tác dụng lên pit-tông cân bằng với lực ma sát giữa pit-tông và xilanh:
			\item Công của khối khí thực hiện có độ lớn: $|A|=Fd=F_{ms}.d=20.0,05=1\ J$.
			\item[b)] Khối khí thực hiện công nên $A=-1\ J$.
			\item Độ biến thiên nội năng của khối khí: $\Delta U=Q+A=1,5-1=0,5\ J$.
		\end{itemize}
	}
\end{vidu}
\begin{vidu}
	Khi truyền một nhiệt lượng Q cho khối khí trong một xilanh hình trụ thì khí dãn nở đẩu pit-tông làm thể tích của khối khí tăng thêm 7 lít. Biết áp suất của khối khí là $3.10^5\ Pa$ và không đổi trong quá trình khí dãn nở. Tính: 
	\begin{itemize}
		\item[a)] Độ lớn của công mà khối khí thực hiện.
		\item[b)] Nhiệt lượng cung cấp cho khối khí. Biết rằng trong quá trình này, nội năng của khối khí giảm 1100 J.
	\end{itemize}
	\loigiai{
		\begin{itemize}
			\item[a)] Độ lớn của công khối khí thực hiện: $|A|=Fd=pSd=p\Delta V=3.10^5.7.10^{-3}=2100\ J$. 
			\item[b)] Do khối khí thực hiện công nên: $A=-2100\ J$
			\item Nội năng của hệ khí giảm 1100 J, nghĩa là độ biến thiên $\Delta U=-1100\ J$.
			\item Áp dụng định luật I nhiệt động lực học, nhiệt lượng cần cung cấp cho khối khí: $Q=\Delta U-A=-1100+2100=1000\ J$.
		\end{itemize}
	}
\end{vidu}
\loai{Tính độ biến thiên nội năng}
\begin{vidu}
	Giả sử cung cấp cho vật một công là 200 J nhưng nhiệt lượng bị thất thoát ra môi trường bên ngoài là 120 J. Hỏi nội năng của vật tăng hay giảm bao nhiêu?
	\loigiai{
		\begin{itemize}
			\item 
			\item 
			\item 
			\item 
			\item 
			\item 
		\end{itemize}
	}
\end{vidu}

\loai{Tính nhiệt lượng}
\begin{vidu}%[1P6-3.2B][Loai2] 
	Người ta thực hiện một công 100 J để nén khí trong xilanh. Biết rằng nội năng của khí tăng thêm 10 J. Khi đó
	\choice
	{khối khí truyền nhiệt là 110 J}
	{khối khí nhận nhiệt là 90 J}
	{khối khí toả nhiệt 110 J}
	{\True khối khí toả nhiệt 90 J}
	\loigiai{
		\Binhluan{
			\begin{itemize}
				\item Thực hiện công lên khối khí suy ra khí nhận công: $A=100\ J$.
				\item Nội năng của khí tăng: $\Delta U=10\ J$.
				\item Nhiệt lượng khối khí: $Q=\Delta U-A=-90\ J.$ \item Khối khí tỏa nhiệt 90 J.
			\end{itemize}
		}
		{$Q<0\Rightarrow$ Khối khí "mất nhiệt", "toả nhiệt", "truyền nhiệt"}
	}
\end{vidu}
\loai{Tính công}
\begin{vidu}%[1P6-3.2B][Loai3]
	Nén một khối khí đựng trong xi lanh với một công A làm khối khí tỏa một nhiệt lượng 40 J. Độ biến thiên nội năng của khối khí là 100 J. Tính A.
	\choice
	{\True 140 J}
	{120 J}
	{70 J}
	{100 J}
	\loigiai{
		\Binhluan{
			\begin{itemize}
				\item Khí tỏa nhiệt: $Q=-40\ J$.
				\item Độ biến thiên nội năng của khí: $\Delta U=100\ J$.
				\item Theo nguyên lý I: 
				\begin{align*}
					\Delta U=Q+A\Rightarrow A=\Delta U-Q=100-(-40)=140\ J.
				\end{align*}	
			\end{itemize}
		}
		{Khi giãn nở khí sinh công. Khi bị nén thì khí nhận công}
	}
\end{vidu}
%\loai{Quá trình đẳng tích}
%\begin{vidu}%[1P6-3.2K][Loai4]
%Cho một bình kín có dung tích coi như không đổi chứa 14 g khí Nitơ ở áp suất 1 atm và nhiệt độ $27^\circ C$. Khí được đun nóng sao cho áp suất tăng gấp 5 lần. Lấy nhiệt dung riêng $c=0,75\ kJ/ (kg.K)$. Nội năng của khí biến thiên một lượng là
%\choice
%{62160 J}
%{\True 12600 J}
%{$-12432$ J}
%{$-62160$ J}
%\loigiai{
	%\Binhluan{
		%\begin{itemize} 
		%\item Vì dung tích không đổi nên V không đổi $\Rightarrow A=0\Rightarrow \Delta U=Q$.
		%\item Độ biến thiên nội năng chính là nhiệt lượng mà khí nhận được. 
		%\item Vì quá trình đẳng tích nên: $p_2=5p_1\Rightarrow T_2=5T_1=1500\ K.$
		%\item Nhiệt lượng mà khí nhận được: 
		%$Q=\Delta U=m.c.\Delta T=12600\ J$.
		%\end{itemize}
		%}
	%{Đổi đơn vị $m=0,014\ kg;\ T=300\ K$ và $c=750\ J/(kg.K)$}
	%}
%\end{vidu}
\loai{Giãn nén khí đẳng áp}
\begin{vidu}%[1P6-3.2K][Loai5] 
	Khi truyền nhiệt lượng $6. 10^6$ J cho khí trong một xilanh hình trụ thì khí nở ra đẩy pit-tông làm thể tích của khí tăng thêm $0,5\ m^3$. Biết áp suất của khí là $8. 10^6\ N/m^2$ và coi áp suất này không đổi trong quá trình khí thực hiện công. Độ biến thiên nội năng của khí là
	\choice
	{$4. 10^6$ J}
	{$6. 10^6$ J}
	{\True $2. 10^6$ J}
	{2 kJ}
	\loigiai{
		\Binhluan{
			\begin{itemize}
				\item Khí nhận nhiệt: $Q=6.10^6\ J$.
				\item Khí dãn nở nên thực hiện công: 
				\begin{align*}
					A=-\left|p.\Delta V\right|=-8.10^6.0,5=-4.10^6\ J.
				\end{align*}
				\item Độ biến thiên nội năng của khí: $\Delta U=A+Q=2.10^6\ J$.
			\end{itemize}
		}
		{Ta biện luận thêm một bước nữa để xác định dấu của A}
	}
\end{vidu}
%\begin{vidu}%[1P6-3.2G][Loai5]
%Một khối khí có $V = 7,5$ lít, $p = 2.10^5\ N/m^2$, $t = 27\doC$ bị nén đẳng áp và nhận một công 50 J. Nhiệt độ của khí sau khi nén là
%\choice
%{$27\doC$}
%{\True $17\doC$}
%{$21\doC$}
%{$25\doC$}
%\loigiai{
	%\begin{itemize}
	%\item Khối khí nhận công, thể tích giảm: $A=p(V_1-V_2)\Rightarrow V_2=V_1-\dfrac{A}{p}=7,5.10^{-3}-\dfrac{50}{2.10^5}=7,25.10^{-3}\ m^3$.
	%\item Áp dụng định luật Gay - Luy sắc: $\dfrac{V_1}{T_1} = \dfrac{V_2}{T_2} \Rightarrow T_2 =T_1.\dfrac{V_2}{V_1}=300.\dfrac{7,25.10^{- 3}}{7,5.10^{-3}}=290\ K=17\doC$.
	%\end{itemize}
	%}
%\end{vidu}
%\loai{Đồ thị}
%\begin{vidu}%[1P6-3.2K][Loai6] 
%immini[thm]{Một khối khí lí tưởng thực hiện quá trình như trên đồ thị. Công mà khối khí trao đổi với môi trường là
	%\choice
	%{0,6 kJ}
	%{\True 0,9 kJ}
	%{1,5 kJ}
	%{1,2 kJ}}{\begin{tikzpicture}[scale=0.8,thick,>=stealth']
		%\draw[->] (0,0)--(4.5,0)node[above]{$V\ (m^3)$}; 
		%\draw[->] (0,0)--(0,3)node[right]{$p\ (kPa)$};
		%\draw[dashed](0,2)node[left]{$300$}--(1,2)--(3,2);
		%\draw[dashed](1,2)--(1,0)node[below=3pt,inner sep=1pt]{\small $0,002$}(3,2)--(3,0)node[below=3pt,inner sep=1pt]{\small $0,005$};
		%\draw[directed,blue, ultra thick](1,2)--(3,2);
		%\end{tikzpicture}}
		%\loigiai{Thể tích tăng, khối khí thực hiện công: 
			%\begin{align*}
			%A=p.(V_2-V_1)=3.10^5.(0,005-0.002)=0,9\ kJ.
			%\end{align*}	
			%}
		%\end{vidu}
		\loai{Tính hiệu suất của động cơ nhiệt}
		\begin{vidu}%[1P6-4.2B][Loai1]  
			Một động cơ nhiệt nhận từ nguồn nóng một nhiệt lượng 1200 J và truyền cho nguồn lạnh một nhiệt lượng 900 J. Hiệu suất của động cơ là
			\choice
			{lớn hơn 75\%}
			{75\%}
			{\True 25\%}
			{nhỏ hơn 25\%}
			\loigiai{
				Hiệu suất động cơ: $H=\dfrac{A}{Q_N}=\dfrac{Q_N-Q_L}{Q_N}=0,25=25\%$
			}
		\end{vidu}
		\loai{Tính nhiệt lượng nguồn nóng của động cơ nhiệt}
		\begin{vidu}%[1P6-4.2B][Loai2] 
			Hiệu suất của động cơ nhiệt là 40\%, động cơ nhiệt thực hiện một công 800 J.  Nhiệt lượng của nguồn nóng là
			\choice
			{480 J}
			{\True 2 kJ}
			{800 J}
			{320 J}
			\loigiai{Nhiệt lượng nhận từ nguồn nóng: $H=\dfrac{A}{Q_N}\Rightarrow Q_N=\dfrac{A}{H}=\dfrac{800}{0,4}=2000\ J=2\ kJ.$
			}
		\end{vidu}
		
		\loai{Tính nhiệt lượng nguồn lạnh của động cơ nhiệt nhận được}
		\begin{vidu}%[1P6-4.2K][Loai4] 
			Một động cơ nhiệt có hiệu suất $25\%$, công suất 30 kW. Nhiệt lượng mà nó tỏa ra cho nguồn lạnh trong 5 giờ làm việc liên tục là
			\choice
			{$16,2. 10^7$ J}
			{\True $162. 10^7$ J}
			{$8,1. 10^7$ J}
			{$81. 10^7$ J}
			\loigiai{ 
				\begin{itemize} 
					\item Công khí thực hiện: $A'=P.t=30.10^3.5.3600=54.10^7\ J$.
					\item Nhiệt lượng truyền cho nguồn lạnh:
					\begin{align*}
						Q'_2=A'\dfrac{(1-H)}{H}=162.10^7\ J.
					\end{align*}
			\end{itemize}}
		\end{vidu}
		
		
		\loai{Bài toán tổng hợp}
		\begin{vidu}%[1P6-4.2G][Loai6] 
			Một động cơ của xe máy có hiệu suất là 20\%. Sau một giờ hoạt động tiêu thụ hết 1 kg xăng có năng suất toả nhiệt là $46.10^6$ J/kg. Công suất của động cơ xe máy là
			\choice
			{\True 2556 W}
			{2946 W}
			{4457 W}
			{5692 W}
			\loigiai{
				\begin{itemize} 
					\item Khi 1 kg xăng cháy hết sẽ tỏa ra nhiệt lượng: 
					\begin{align*}
						Q=m.q=46.10^6\ J \Rightarrow
						H=\dfrac{\left|A\right|}{Q}=0,2\Rightarrow \left|A\right|=0,2.Q=0,2.46.10^6=92.10^5\ J.
					\end{align*}
					\item Công suất của động cơ xe máy là: $P=\dfrac{A}{t}=\dfrac{{92.10}^5}{3600}=2556\ W$.
				\end{itemize}
			}
		\end{vidu}
